Las funciones de pérdida son parte crucial del aprendizaje automático. El término refiere a una función matemática que estima el error entre la estimación y el objetivo. En función de la función de pérdida elegida, se prioriza un tipo de error u otro. Por ejemplo:

\begin{itemize}
    \item Función de pérdida L1, error absoluto medio: minimiza el error absoluto promedio entre predicción y objetivo.
\begin{center}
    \[
\mathcal{L}_{L1}(\hat{Y}, Y) =
\frac{1}{N}
\sum_{i=1}^{N}
\left| \hat{Y}_i - Y_i \right|
\]
\end{center}

Donde \(Y_i\) representa el valor real de la magnitud de la fuente en una celda tiempo-frecuencia, \(\hat{Y}_i\) representa el valor estimado por el modelo y \(N\) es el número total de elementos comparados.

    \item Función de pérdida L2, error cuadrático medio: minimiza el error cuadrático medio.
\begin{center}
    \[
\mathcal{L}_{L2}(\hat{Y}, Y) =
\frac{1}{N}
\sum_{i=1}^{N}
\left( \hat{Y}_i - Y_i \right)^2
\]
\end{center}

Donde \(Y_i\) representa el valor real de la magnitud de la fuente, \(\hat{Y}_i\) representa el valor estimado por el modelo y \(N\) es el número total de elementos comparados.


\end{itemize}